\documentclass[11pt]{article}
\usepackage[margin=1in]{geometry}
\usepackage{amsmath,amssymb,amsthm}
\usepackage[T1]{fontenc}
\usepackage{hyperref}
\newcommand{\NE}{\operatorname{NE}}
\newcommand{\Stab}{\operatorname{Stab}_0}
\newcommand{\Reg}{\operatorname{Reg}}
\newcommand{\SwapReg}{\operatorname{SwapReg}}
\newcommand{\pos}[1]{[#1]_+}
\newtheorem{proposition}{Proposition}
\title{Axiomatic containment of Mertens-stable sets without genericity}
\author{}
\date{Status checked: 17 October 2026}
\begin{document}
\maketitle
\section*{Status and provenance}
This is a proposed precise generalization of a literature theorem. Its
formulation is editorial, rather than a verbatim named conjecture.
No solution or proof is claimed. The primary source discusses the
nongeneric gap; it does not establish the proposed universal statement.
\section*{Objects and conventions}
A finite normal-form game $G=(I,(A_i),(u_i))$ has mixed-strategy space
$\Sigma=\prod_{i\in I}\Delta(A_i)$ and Nash set $\NE(G)$.
A solution concept is a rule $\Phi$ defined on \emph{all} finite
normal-form games, including games with additional players and actions:
\[
 \varnothing\ne\Phi(G)\subseteq
 \{K\subseteq\NE(G):K\ne\varnothing,\ K\text{ closed and connected}\}.
\]
Write $\Stab(G)$ for Mertens $0$-stable sets, with rational cohomology,
using exactly Definition 2.4 of the version of Govindan--Wilson cited
below. The zero specifies the coefficient characteristic, not an
equilibrium's local index. To fix the perturbation model, set
\[
 P_\varepsilon=\{\delta\tau:0\leq\delta\leq\varepsilon,\ \tau\in\Sigma\},
 \quad \Sigma_i(\eta)=\{x_i\in\Delta(A_i):x_i(a)\geq\eta_i(a)\},
\]
and let $E_G$ be the equilibrium graph of these constrained games,
projected to $P_1$. The definition combines cohomological essentiality
of this projection with the source's connexity condition. Definition 2.4,
including its boundary conventions, is normative here.

Use the exact cross-game axioms $B,D^*,I^*$ of that same version,
Sections 3--5 (Definitions 4.4 and 5.2), rather than informal substitutes.
Axiom $B$ requires every selected set to contain a profile equivalent to
a quasi-perfect equilibrium in every perfect-recall extensive-form
representation. Axiom $D^*$ concerns additions of strongly irrelevant
dominated strategies. Axiom $I^*$ concerns strong embeddings, including
additional players; its comparisons use the source's projection maps.
The full technical restrictions in those definitions are part of the
problem. Arbitrary dominated-action additions or arbitrary embeddings
are not being substituted for them.

\section*{Question S1}
Suppose $\Phi$ satisfies $B,D^*,I^*$ globally. Is the following true?
\[
 \forall G\quad\forall K\in\Phi(G)\quad
 \exists S\in\Stab(G)\qquad S\subseteq K.
\]
There is no genericity, regularity, or isolated-equilibrium assumption.
A negative answer must specify a globally defined solution concept
satisfying the cross-game axioms, and a finite game with a selected set
that contains no stable set. A bad set in a single game is insufficient.

\section*{Known boundary and motivation}
Govindan--Wilson prove the containment conclusion for generic
perfect-recall extensive-form games \cite[Theorem 6.1]{gw}. Their
Sections 6.1 and 10 separate containment from selection of stable sets
themselves and identify an obstacle in the nongeneric case. The question
above asks whether these same axioms already suffice, rather than
assuming that the generic proof extends.

\section*{Required distinctions}
The three statements
\[
 (\forall K\in\Phi(G)\ \exists S\in\Stab(G):S\subseteq K),\quad
 \Phi(G)\subseteq\Stab(G),\quad
 \Phi(G)=\Stab(G)
\]
have different logical strength. S1 asks only the first. Its proof must
also cover equilibrium continua. It neither imposes stable-set membership
as an axiom nor silently replaces quasi-perfection with sequential
equilibrium in nongeneric representations.

\section*{Checks and remaining uncertainty}
The generic case is covered by the cited theorem. The general
formulation preserves the source's all-game domain for the axioms and
allows the target game itself to be nongeneric. No counterexample search
or new proof was performed. A catalogue and source search found no exact
existing we-solve target. Open status and novelty of this precise
editorial formulation remain subject to expert review.

\section*{Preparation}
The original research map was generated using OpenAI GPT 6 Astra Pro.
Source checks, formulation, and LaTeX preparation used a Codex assistant
based on GPT-6. Exact model versions and reasoning settings are unknown.
The cited researchers retain credit for their results. Source statements
and scope were checked; no complete proof audit or novelty certification
is claimed.
\paragraph{Origin of the questions.}
The open-problem suggestions motivating this entry are credited to
Rida Laraki's Facebook comment~\cite{laraki-facebook}.
\begin{thebibliography}{9}
\bibitem{gw}
Srihari Govindan and Robert Wilson.
\emph{Axiomatic Equilibrium Selection in Generic Extensive-Form Games}.
arXiv:2504.16908v2, 26 August 2026.
Definitions 2.4, 2.5, 4.4, 5.2; Axioms 3.1, 4.7, 5.4; Theorem 6.1; Sections 6.1 and 10.
\url{https://arxiv.org/html/2504.16908v2}.
\bibitem{laraki-facebook}
Rida Laraki.
\emph{Facebook comment suggesting open problems in game theory}.
\end{thebibliography}
\end{document}
